\section{Un corrector que usa sólo el punto estacionario de torres}
Sea $A$ una máscara binaria de tamaño $k$. En su soporte existe un punto estacionario
interior $P=(p_{ij})$ del problema de torres, bistocástico, con $p_{ij}\le M/k$.
Ponemos $J=\mathbf1\mathbf1^\top/k$, $B=P-J$ y suponemos $\|B\|_{\rm op}\le r<1$.
Todas las constantes de los restos dependen sólo de $M,r$. Suponemos $k$ suficientemente grande.
Las sumas por sitio recorren las $k^2$ posiciones. Una potencia $\circ m$ actúa por entradas.
Definimos
\[
 G_L=BB^\top,\quad G_R=B^\top B,\quad
 T_s=G_s(I-G_s)^{-1}\quad(s=L,R),\qquad U=BT_R .
\]
Sea $j$ el vector de longitud $2k$ que concatena las diagonales de $T_s-G_s$.
Sean $j_0=\tr(T_R-G_R)$, $q_4=\tr G_R^2$,
$\Sigma_3(X)=\sum_{ij}X_{ij}^3$ y
\[
 H=\begin{pmatrix}-T_L&B+U\\ (B+U)^\top&-T_R\end{pmatrix}.
\]
\paragraph{Fórmula de primer orden.}
Para $g_T(A)=\log\per A-L_T(A)$, con la normalización de torres de C,
\begin{align}
 g_T(A)&=-1/(12k)+\Xi(P)+O(k^{-2}),\\
 \Xi(P)&=\frac{j_0^2/2+\tr T_R^2-3j_0/2}{k}
       -\tfrac34\|j\|_2^2+\tfrac12 j^\top Hj \notag\\
 &\quad-\tfrac13\{\Sigma_3(T_L)+\Sigma_3(T_R)\}
       +2\sum_{ij}B_{ij}U_{ij}^2+\tfrac23\Sigma_3(U),\label{eq:xi}\\
 |\Xi(P)|&\le C(M,r)q_4/k .\label{eq:q}
\end{align}
La matriz $P$ puede tener ceros estructurales y grados desiguales.
Este resultado utiliza el desarrollo del permanente con resto uniforme $O(k^{-2})$
admitido en la ronda 1. Las dos secciones siguientes dan la deducción.
\paragraph{Qué mide $q_4$.}
Como $PP^\top-J=BB^\top$, se tiene
\[
 q_4=\sum_{ij}\big((PP^\top)_{ij}-1/k\big)^2
 =\sum_\ell \sigma_\ell(B)^4 .
\]
Así se puede evaluar el exceso ponderado de ciclos de longitud cuatro con repeticiones permitidas.
Si $I,J'$ ocupan fracciones $a,b\in(0,1)$ de las filas y columnas y
$d=|\sum_{I\times J'}B_{ij}|/k$, entonces
$q_4\ge d^4/[a(1-a)b(1-b)]^2$.
En efecto, los indicadores centrados de ambos conjuntos tienen normas $\sqrt{ka(1-a)}$
y $\sqrt{kb(1-b)}$. Su producto bilineal con $B$ tiene módulo $kd$.
Esta cota inferior sobre $q_4$ no fija el signo de $\Xi$.
